\documentclass[10pt,a4paper]{article}
\usepackage[utf8]{inputenc}
\usepackage[a4paper,left=1.35cm,right=1.35cm,bottom=1.35cm,top=2cm]{geometry}
\usepackage{xcolor}
\usepackage[hidelinks]{hyperref}
\usepackage{enumitem}
\usepackage{titlesec}

\pagestyle{empty}
\flushbottom
\setlength{\parindent}{0pt}
\setlength{\parskip}{3pt plus 1pt minus 1pt}
\setlist[itemize]{leftmargin=*,itemsep=2pt,topsep=3pt,parsep=0pt,partopsep=0pt}
\setcounter{secnumdepth}{0}

\titleformat{\section}{\large\bfseries\color{blue!65!black}}{}{0pt}{}[{\color{blue!65!black}\titlerule}]
\titlespacing{\section}{0pt}{6pt plus 1pt minus 1pt}{3pt plus 1pt minus 1pt}

\begin{document}

{\centering
{\LARGE \textbf{Amira Boubaker}}\\[7pt]
Développeuse Full-Stack Junior \\[4pt]
Monastir, Tunisia \quad|\quad +216-24-503-203 \quad|\quad amiraboubakeresprims@gmail.com\\[4pt]
\href{https://amiraboubaker.github.io/AmiraBoubaker-portfolio/}{Portfolio} \quad|\quad
\href{https://www.linkedin.com/in/amira-boubaker/}{LinkedIn} \quad|\quad
\href{https://github.com/amiraboubaker}{GitHub} \quad|\quad
\href{https://leetcode.com/u/Amira2023/}{LeetCode}
\par}

\section{Profil}
Développeuse full-stack orientée métier, spécialisée dans la conception d'applications web robustes et maintenables. Expérience pratique en JavaScript/TypeScript, Node.js, SQL, API REST et intégration de fonctionnalités utiles au quotidien des utilisateurs. Approche analytique pour transformer un besoin fonctionnel en solution technique claire.

\section{Compétences Clés}
\textbf{Front-end :} HTML, CSS, JavaScript, ReactJS, React Native, TypeScript

\textbf{Back-end :} Node.js, Express.js, Symfony 6, .NET, API REST, POO

\textbf{Données \& intégration :} MySQL, MongoDB, Firebase, Oracle, SQL Server, intégration de services

\textbf{Méthodes \& outils :} Agile/Scrum, Git, UML, VS Code, travail en équipe pluridisciplinaire

\section{Expérience}
\textbf{Software Engineer, ESPITA (Sousse)} \hfill \textit{Mar 2026 -- Present}
\begin{itemize}
  \item Développement et livraison de trois applications web institutionnelles pour ESPITA, ESEAC Sfax et EMP Sousse.
  \item Conception de chatbots orientés usage pour améliorer l'accès à l'information et l'autonomie des utilisateurs.
  \item Participation aux choix techniques et à la qualité du code pour garantir maintenabilité, fiabilité et évolutivité.
\end{itemize}

\textbf{React Native Developer, Mobelite (Monastir)} \hfill \textit{Feb 2025 -- Jul 2025}
\begin{itemize}
  \item Livraison de JourneyBuddy, une application de voyage assistée par l'IA utilisée par plus de 50 000 utilisateurs.
  \item Intégration de services externes (cartographie, assistant conversationnel) pour enrichir les parcours fonctionnels.
  \item Mise en place d'une architecture modulaire pour faciliter l'évolution produit et la maintenance.
  \item Gestion du cycle de publication mobile, de la génération à la distribution des APK/AAB.
\end{itemize}

\textbf{Software Engineer Student, ESPRIM (Monastir)} \hfill \textit{Oct 2024 -- May 2025}
\begin{itemize}
  \item Développement d'une plateforme e-commerce full-stack (Ktabinet) pour la vente de livres en ligne.
  \item Implémentation de workflows métier (stock, paiement) avec Symfony 6 et MySQL.
  \item Traduction de besoins de gestion en fonctionnalités concrètes via API, logique métier et modèle de données relationnel.
\end{itemize}

\textbf{Software Engineer Intern, Nexans (Monastir)} \hfill \textit{Jul 2024 -- Sep 2024}
\begin{itemize}
  \item Développement d'ECTRH, application RH mobile pour la gestion des données collaborateurs et le suivi des équipes.
  \item Automatisation de processus pour fluidifier les opérations et simplifier la gestion quotidienne.
  \item Collaboration étroite avec les parties fonctionnelles et techniques pour livrer des solutions alignées au besoin métier.
\end{itemize}

\section{Projets}

\textbf{Eseac (2026)} \hfill \textit{Application web full-stack avec chatbot et tableau d'administration}
\begin{itemize}
  \item Application web orientée service, permettant aux étudiants, candidats et visiteurs d'interagir avec les contenus et actualités.
  \item Conception d'un espace d'administration pour la gestion structurée des utilisateurs et des informations.
  \item Mise en pratique des logiques de flux, de données et de processus proches des besoins métier.
  \item GitHub: \href{https://github.com/amiraboubaker/eseac/}{github.com/amiraboubaker/eseac/}
\end{itemize}

\textbf{JourneyBuddy (2025)} \hfill \textit{React Native, Node.js, MongoDB, AI Chatbot}
\begin{itemize}
  \item Application de planification de voyage avec recommandations personnalisées.
  \item Démonstration de l'intégration API et de la gestion de données applicatives à grande échelle d'usage.
  \item GitHub: \href{https://github.com/dghama/journeyBuddy-INT2}{github.com/dghama/journeyBuddy-INT2}
\end{itemize}

\textbf{HealthTracker (2025)} \hfill \textit{React Native, Firebase, AI Chatbot}
\begin{itemize}
  \item Application bien-être pour le suivi du sommeil, de la nutrition et de l'hydratation avec coaching IA.
  \item Renforcement des compétences en interfaces interactives, structuration de données et expérience utilisateur.
  \item GitHub: \href{https://github.com/amiraboubaker/HealthTracker}{github.com/amiraboubaker/HealthTracker}
\end{itemize}

\section{Formation \& Certification}
\textbf{Diplôme d'Ingénierie en Informatique}, ESPRIM (Monastir) \hfill \textit{Sep 2022 -- Jul 2025}\\
Spécialisation : Génie Logiciel

\textbf{AI Engineer for Data Scientists -- Associate}, DataCamp

\end{document}
